\documentclass[12pt,a4paper]{article}

\usepackage[margin=2.5cm]{geometry}
\usepackage{hyperref}
\usepackage{listings}
\usepackage{xcolor}
\usepackage{booktabs}
\usepackage{array}
\usepackage{enumitem}
\usepackage{amsmath}
\usepackage{amssymb}
\usepackage{titlesec}
\usepackage{parskip}
\usepackage{fancyhdr}
\usepackage{graphicx}

\hypersetup{
    colorlinks=true,
    linkcolor=blue!60!black,
    urlcolor=blue!60!black,
}

\lstdefinestyle{cpp}{
  language=C++,
  basicstyle=\ttfamily\small,
  keywordstyle=\color{blue!70!black}\bfseries,
  commentstyle=\color{green!50!black},
  stringstyle=\color{red!60!black},
  numbers=left,
  numberstyle=\tiny\color{gray},
  numbersep=6pt,
  breaklines=true,
  frame=single,
  rulecolor=\color{gray!40},
  backgroundcolor=\color{gray!5},
  tabsize=4,
  showstringspaces=false,
  xleftmargin=1.5em,
}

\lstset{style=cpp}

\pagestyle{fancy}
\fancyhf{}
\rhead{COP290 -- Assignment 4}
\lhead{DataFrameLib Design Report}
\cfoot{\thepage}

\title{%
  \Large\textbf{DataFrameLib} \\[0.4em]
  \large Design Report \\[0.2em]
  \normalsize COP290 -- Design Practices
}
\author{Amar Sinha}
\date{April 2026}

\begin{document}

\maketitle
\tableofcontents
\newpage

%--------------------------------------------------------------------------
\section{Overview}
%--------------------------------------------------------------------------

DataFrameLib is a C++ library for high-performance tabular data processing,
modelled after the ideas behind Pandas and Polars.  It provides two
execution modes and a query optimiser:

\begin{itemize}[leftmargin=1.5em]
  \item \textbf{EagerDataFrame} -- immediate, materialised execution.
  \item \textbf{LazyDataFrame} -- deferred execution via a computation DAG.
  \item \textbf{QueryOptimizer} -- rule-based plan optimisation before
        materialisation.
\end{itemize}

All in-memory columnar storage and I/O (CSV and Parquet) is backed
exclusively by \textbf{Apache Arrow}.  The build system is CMake (C++20)
with shared-library linkage against \texttt{Arrow::arrow\_shared} and
\texttt{Parquet::parquet\_shared}.

%--------------------------------------------------------------------------
\section{Project Structure}
%--------------------------------------------------------------------------

\begin{verbatim}
.
├── CMakelists.txt
├── include/
│   └── dataframelib/
│       ├── types.hpp          -- DType enum + Arrow type utilities
│       ├── column.hpp         -- Column value type + column-level ops
│       ├── expr.hpp           -- Expression DSL (Expr, ExprNode hierarchy)
│       ├── plan_node.hpp      -- Logical plan nodes (PlanNode hierarchy)
│       ├── eager_dataframe.hpp
│       ├── lazy_dataframe.hpp
│       └── optimizer.hpp
└── src/
    ├── column.cpp             -- Arithmetic/comparison/string/null kernels
    ├── expr.cpp               -- Expr factory functions
    ├── plan_node.cpp          -- ToString() + BuildDotGraph()
    ├── eager_dataframe.cpp    -- Eager execution engine
    ├── lazy_dataframe.cpp     -- Lazy plan builder + collect() executor
    └── optimizer.cpp          -- QueryOptimizer + all rule passes
\end{verbatim}

%--------------------------------------------------------------------------
\section{Type System}
%--------------------------------------------------------------------------

Supported column types are declared in \texttt{types.hpp}:

\begin{center}
\begin{tabular}{ll}
\toprule
\textbf{DType enum} & \textbf{Arrow type} \\
\midrule
\texttt{kInt32}   & \texttt{arrow::int32()} \\
\texttt{kInt64}   & \texttt{arrow::int64()} \\
\texttt{kFloat32} & \texttt{arrow::float32()} \\
\texttt{kFloat64} & \texttt{arrow::float64()} \\
\texttt{kString}  & \texttt{arrow::utf8()} \\
\texttt{kBoolean} & \texttt{arrow::boolean()} \\
\bottomrule
\end{tabular}
\end{center}

\textbf{Type promotion rule.}  When a binary numeric operation mixes
integer and floating-point operands the result type is resolved by
\texttt{ResolvePromotedNumericType()} as: \texttt{double} if either side
is \texttt{double}; otherwise \texttt{float} if either side is
\texttt{float}; otherwise \texttt{int64} if either side is \texttt{int64};
otherwise \texttt{int32}.  This covers the required rule $\mathtt{int +
float = float}$.

\textbf{Missing values.}  The \texttt{Column::Make()} factory calls
\texttt{RejectNaNInFloatingArrays()}, which scans floating-point arrays on
construction and raises \texttt{arrow::Status::Invalid} if any NaN is
found.  Missing values must be represented as Arrow nulls only.

Incompatible operations (e.g.\ arithmetic on a string column) throw
\texttt{arrow::Status::TypeError} immediately.

%--------------------------------------------------------------------------
\section{Column Kernel Layer}
%--------------------------------------------------------------------------

\texttt{column.cpp} / \texttt{column.hpp} define a set of free functions
that operate on \texttt{Column} values:

\begin{itemize}[leftmargin=1.5em]
  \item \textbf{Binary numeric:} \texttt{ApplyBinaryNumericOp} --
        addition, subtraction, multiplication, division, modulo with
        type-promotion and null propagation.
  \item \textbf{Unary absolute value:} \texttt{ApplyUnaryAbsOp} --
        dispatches per numeric type.
  \item \textbf{Comparison:} \texttt{ApplyComparisonOp} -- supports all
        six relational operators across numeric, string and boolean columns.
        Any null operand produces a null result.
  \item \textbf{Boolean:} \texttt{ApplyBooleanBinaryOp},
        \texttt{ApplyBooleanNotOp} -- AND, OR, NOT on boolean columns.
  \item \textbf{Null predicates:} \texttt{ApplyNullPredicateOp} --
        \texttt{is\_null} / \texttt{is\_not\_null}.
  \item \textbf{String unary:} \texttt{ApplyStringUnaryOp} -- length
        (returns \texttt{int64}), to\_lower, to\_upper.
  \item \textbf{String predicates:} \texttt{ApplyStringPredicateOp} --
        contains, starts\_with, ends\_with (return \texttt{boolean}).
\end{itemize}

All kernels are hand-written (no Arrow compute kernel dependency) and
propagate nulls element-wise.

%--------------------------------------------------------------------------
\section{Expression System}
%--------------------------------------------------------------------------

The expression DSL lives in \texttt{expr.hpp} / \texttt{expr.cpp}.  An
\texttt{Expr} is a thin wrapper around a \texttt{shared\_ptr<ExprNode>}.
Operator overloads build the expression tree without evaluating it,
enabling the same \texttt{Expr} objects to be used in both eager and lazy
contexts.

\subsection{ExprNode Hierarchy}

\begin{center}
\begin{tabular}{lll}
\toprule
\textbf{Node class} & \textbf{Kind} & \textbf{Description} \\
\midrule
\texttt{ColNode}      & Col      & Reference to a named column \\
\texttt{LitNode}      & Lit      & Scalar literal (Arrow Scalar) \\
\texttt{AliasNode}    & Alias    & Renames the output of a sub-expression \\
\texttt{BinaryNode}   & Binary   & Two-child arithmetic/comparison/boolean node \\
\texttt{UnaryNode}    & Unary    & Single-child (not, abs, is\_null, is\_not\_null) \\
\texttt{StringOpNode} & StringOp & String function with optional string argument \\
\texttt{AggNode}      & Agg      & Aggregation function (sum, mean, count, min, max) \\
\bottomrule
\end{tabular}
\end{center}

\subsection{Free-function Factories}

\begin{lstlisting}
Expr col(std::string name);
Expr lit(int32_t value);   Expr lit(int64_t value);
Expr lit(float value);     Expr lit(double value);
Expr lit(bool value);      Expr lit(std::string value);
\end{lstlisting}

\subsection{Operator Overloads on Expr}

\begin{center}
\begin{tabular}{ll}
\toprule
\textbf{C++ syntax} & \textbf{Semantic} \\
\midrule
\texttt{e1 + e2} / \texttt{-} / \texttt{*} / \texttt{/} / \texttt{\%}
    & Arithmetic (BinaryNode) \\
\texttt{e1 == e2} / \texttt{!=} / \texttt{<} / \texttt{<=} / \texttt{>} / \texttt{>=}
    & Comparison (BinaryNode) \\
\texttt{e1 \& e2} / \texttt{e1 | e2} / \texttt{\textasciitilde{}e}
    & Boolean AND / OR / NOT \\
\texttt{e.abs()} & Absolute value \\
\texttt{e.is\_null()} / \texttt{e.is\_not\_null()} & Null predicates \\
\texttt{e.length()} / \texttt{e.to\_lower()} / \texttt{e.to\_upper()} & String ops \\
\texttt{e.contains(s)} / \texttt{e.starts\_with(s)} / \texttt{e.ends\_with(s)}
    & String predicates \\
\texttt{e.sum()} / \texttt{e.mean()} / \texttt{e.count()} / \texttt{e.min()} / \texttt{e.max()}
    & Aggregations \\
\texttt{e.alias(name)} & Output rename \\
\bottomrule
\end{tabular}
\end{center}

%--------------------------------------------------------------------------
\section{EagerDataFrame}
%--------------------------------------------------------------------------

\texttt{EagerDataFrame} stores data as a
\texttt{std::vector<Column>} (in declaration order) plus an
\texttt{arrow::Schema} and a row count.  Every method materialises its
result immediately and returns a new \texttt{EagerDataFrame} by value --
no in-place mutation.

\subsection{I/O}

\begin{itemize}[leftmargin=1.5em]
  \item \textbf{read\_csv} -- uses \texttt{arrow::csv::TableReader};
        chunked arrays are flattened to contiguous arrays via
        \texttt{arrow::Concatenate}.
  \item \textbf{read\_parquet} -- uses
        \texttt{parquet::arrow::OpenFile} + \texttt{ReadTable}.
  \item \textbf{write\_csv} / \textbf{write\_parquet} -- builds a
        temporary \texttt{arrow::Table} and delegates to
        \texttt{arrow::csv::WriteCSV} / \texttt{parquet::arrow::WriteTable}.
  \item \textbf{from\_columns(map)} -- accepts an
        \texttt{unordered\_map<string, shared\_ptr<Array>>}; validates equal
        lengths; sorts keys alphabetically for deterministic column order.
\end{itemize}

\subsection{Transformations}

\begin{description}
  \item[select(columns)] Picks a subset by name, preserving the requested
        order.  Duplicate names throw immediately.

  \item[select(expressions)] Evaluates each \texttt{Expr} via
        \texttt{EvaluateNodeToColumn}, which recursively dispatches on
        \texttt{ExprNode::Kind}.  The output name defaults to the column
        reference name, the alias name if an \texttt{AliasNode} wraps the
        expression, or \texttt{expr\_<i>} otherwise.

  \item[filter(predicate)] Evaluates the predicate expression to a
        boolean column, then passes through \texttt{FilterArrayByMask}
        for each column.  Null mask entries are treated as \texttt{false}
        (row excluded).

  \item[with\_column(name, expr)] Evaluates \texttt{expr}; if a column
        with the same name exists it is replaced in place; otherwise the
        new column is appended.

  \item[group\_by(keys).aggregate(aggs)] The group-by is implemented with
        a hash-based approach: \texttt{CellToKeyString} serialises each
        key column cell into a type-tagged string; groups are built into an
        \texttt{unordered\_map} that preserves insertion order via a
        separate \texttt{group\_order} vector.
        \texttt{AggregateArrayByGroups} then computes sum, mean, count,
        min, or max per group, returning \texttt{null} when all values in
        the group are null.

  \item[join(other, on, how)] Supports inner, left, right, and outer joins.
        A hash index is built on the right table; the left table is probed.
        Null key values on either side are excluded from matching.
        Non-key columns from both sides are appended; duplicate non-key
        column names raise an error before execution begins.

  \item[sort(columns, ascending)] A stable sort on a permutation index
        vector using a custom comparator that handles multi-key tie-breaking
        and null-last ordering.

  \item[head(n)] Slices each column array via \texttt{arrow::Array::Slice}.
\end{description}

%--------------------------------------------------------------------------
\section{LazyDataFrame}
%--------------------------------------------------------------------------

\texttt{LazyDataFrame} holds a single \texttt{shared\_ptr<PlanNode>}.
Every transformation method creates a new plan node that wraps the current
plan and returns a new \texttt{LazyDataFrame} by value.  No data is
read or processed until \texttt{collect()} is called.

\subsection{Plan Node Hierarchy}

\begin{center}
\begin{tabular}{ll}
\toprule
\textbf{Node class} & \textbf{Logical operation} \\
\midrule
\texttt{ScanCsvNode}           & Leaf -- CSV file path \\
\texttt{ScanParquetNode}       & Leaf -- Parquet file path \\
\texttt{SelectNamesNode}       & Column projection by name \\
\texttt{SelectExprsNode}       & Column projection by expression \\
\texttt{FilterNode}            & Row predicate \\
\texttt{WithColumnNode}        & Derived column addition / replacement \\
\texttt{GroupByAggregateNode}  & Combined group-by + aggregation \\
\texttt{JoinNode}              & Binary join (left + right inputs) \\
\texttt{SortNode}              & Row ordering \\
\texttt{HeadNode}              & Row-count limit \\
\bottomrule
\end{tabular}
\end{center}

All nodes except \texttt{JoinNode} have a single primary \texttt{input()}.
\texttt{JoinNode} additionally exposes \texttt{right\_input()}, making the
plan a DAG rather than a simple tree in the presence of joins.

\subsection{Execution via collect()}

\texttt{collect()} runs the \texttt{QueryOptimizer} on the plan and then
calls \texttt{ExecutePlan(optimized\_plan)}.  The executor performs a
post-order traversal of the plan DAG using an explicit stack (to avoid
stack-overflow on deep plans), materialises each node into an
\texttt{EagerDataFrame}, and stores results in an
\texttt{unordered\_map<const PlanNode*, EagerDataFrame>}.  Each node
retrieves its input(s) from this map and calls the corresponding
\texttt{EagerDataFrame} method.

\subsection{sink\_csv / sink\_parquet}

These simply call \texttt{collect()} and forward to
\texttt{EagerDataFrame::write\_csv/write\_parquet}.

\subsection{explain(path)}

\texttt{explain(path)} renders \emph{two} PNG diagrams side-by-side:

\begin{enumerate}
  \item \texttt{<base>\_logical.png} -- the unoptimised plan DAG.
  \item \texttt{<base>\_optimized.png} -- the plan after
        \texttt{QueryOptimizer::optimize()}.
\end{enumerate}

The DOT source is generated by \texttt{BuildDotGraph()} in
\texttt{plan\_node.cpp}, which assigns integer node IDs, emits box-shaped
nodes labelled with \texttt{PlanNode::ToString()}, and draws directed
edges from child to parent (data flows upward).  The DOT file is then
rendered to PNG by shelling out to \texttt{dot -Tpng} from Graphviz; the
temporary \texttt{.dot} file is cleaned up afterwards.

%--------------------------------------------------------------------------
\section{Query Optimiser}
%--------------------------------------------------------------------------

\texttt{QueryOptimizer::optimize(root)} applies rules in multiple passes
using \texttt{OptimizeTreeOnce()}, which traverses the plan bottom-up,
rebuilds each node with optimised children, and applies all rule functions
to the rebuilt node.  Convergence is detected by a string fingerprint of
the plan tree; the loop runs for at most 8 passes.

The following optimisations are implemented.

%------------------------------------------------------------------
\subsection{Constant Folding \& Expression Simplification}
%------------------------------------------------------------------

\subsubsection*{Transformation}

\texttt{OptimizeExpr(expr)} walks an expression tree bottom-up.  If both
children of a binary node or the child of a unary/string node reduce to
literal values, the expression is evaluated at plan-construction time and
replaced by a single \texttt{LitNode} carrying the result.

Additionally, the following algebraic identities are applied even when only
one operand is a literal:

\begin{itemize}[leftmargin=1.5em]
  \item $x + 0 \;\to\; x$ \quad and \quad $0 + x \;\to\; x$
  \item $x - 0 \;\to\; x$
  \item $x \times 1 \;\to\; x$ \quad and \quad $1 \times x \;\to\; x$
  \item $x / 1 \;\to\; x$
  \item $\mathtt{true} \;\&\; e \;\to\; e$ \quad and \quad $e \;\&\; \mathtt{true} \;\to\; e$
  \item $\mathtt{false} \;|\; e \;\to\; e$ \quad and \quad $e \;|\; \mathtt{false} \;\to\; e$
  \item $\sim(\sim e) \;\to\; e$ (double negation elimination)
\end{itemize}

String and unary literal folds are also performed: e.g.\
\texttt{lit("hello").length()} $\to$ \texttt{lit(5)}, and
\texttt{lit(true).is\_null()} $\to$ \texttt{lit(false)}.

\subsubsection*{Correctness}

All identities are standard algebraic equalities that hold over the
nullable semantics defined for DataFrameLib.  Division-by-zero and
modulo-by-zero are guarded: the fold is skipped (returning the original
expression) if the divisor is zero.  The same applies to type mismatches
(e.g.\ modulo on floating-point literals).

\subsubsection*{Example}

\begin{lstlisting}
col("salary") * lit(1) + lit(0)
// After constant folding + simplification:
col("salary")
\end{lstlisting}

\subsubsection*{Performance benefit}

Eliminates redundant arithmetic and reduces the number of expression nodes
that the execution engine must visit at runtime.  For wide tables this
avoids allocating and populating temporary intermediate arrays.

%------------------------------------------------------------------
\subsection{Constant Filter Elimination}
%------------------------------------------------------------------

\subsubsection*{Transformation}

After \texttt{MakePlanWithOptimizedExprs} folds all literal sub-expressions,
\texttt{ApplyConstantFilterElimination} inspects the predicate of a
\texttt{FilterNode}.  If the predicate is a boolean literal it is handled as
follows:

\begin{itemize}[leftmargin=1.5em]
  \item \textbf{Predicate is \texttt{lit(true)}:} The filter node is removed
        entirely and replaced by its child -- no row is ever excluded so the
        node is pure identity overhead.
  \item \textbf{Predicate is \texttt{lit(false)}:} The filter is replaced by
        \texttt{HeadNode(input, 0)}.  This canonicalises ``zero rows
        possible'' into the existing \texttt{head} mechanism so that
        subsequent limit pushdown passes can propagate the zero-row signal
        as close to the scan as possible.
\end{itemize}

\subsubsection*{Correctness}

$\sigma_{\mathtt{true}}(R) = R$ and $\sigma_{\mathtt{false}}(R) = \emptyset$
are set-algebra identities.  The \texttt{false} rewrite is safe because
\texttt{head(0)} on any relation always produces zero rows in the same
schema, matching the semantics of an impossible filter.  The pass runs
immediately after constant folding so that predicates like
\texttt{lit(1)==lit(1)} are already reduced to \texttt{lit(true)} before
this rule fires.

\subsubsection*{Example}

\begin{lstlisting}
// Original (tautology predicate that folds to true)
scan_csv("employees.csv")
  .filter(lit(1) == lit(1))
  .select(...)

// After constant folding: lit(1)==lit(1) -> lit(true)
// After constant filter elimination: filter node removed
scan_csv("employees.csv")
  .select(...)
\end{lstlisting}

\subsubsection*{Performance benefit}

Removes an entire filter evaluation pass -- no boolean array is allocated,
no mask scan is performed.  For the \texttt{lit(false)} case, combined with
limit pushdown it can eventually prevent reading the source file at all.

%------------------------------------------------------------------
\subsection{Predicate Pushdown}
%------------------------------------------------------------------

\subsubsection*{Transformation}

\texttt{ApplyPredicatePushdown(node)} inspects a \texttt{FilterNode} and
attempts to push it closer to the data source by examining the node
immediately below it:

\begin{itemize}[leftmargin=1.5em]
  \item \textbf{Filter over SelectNames:} If all columns referenced by the
        predicate are present in the projected set, the filter is pushed
        below the projection: \texttt{select(cols).filter(pred)} $\to$
        \texttt{filter(pred).select(cols)}.

  \item \textbf{Filter over WithColumn:} If the predicate does not
        reference the newly derived column, the filter is pushed below the
        \texttt{with\_column}: \texttt{with\_column(c,e).filter(pred)} $\to$
        \texttt{filter(pred).with\_column(c,e)}.

  \item \textbf{Filter over Filter:} Two consecutive filters are merged
        into a single filter whose predicate is the AND of both.  The merged
        predicate is itself passed through \texttt{OptimizeExpr}.

  \item \textbf{Filter over Sort:} The filter is pushed below the sort
        (sorting fewer rows is cheaper): \texttt{sort(...).filter(pred)}
        $\to$ \texttt{filter(pred).sort(...)}.

  \item \textbf{Filter over Join (inner / left / right):} If the predicate
        references only join-key columns, it is pushed into the appropriate
        side(s): both sides for inner joins, only the left for left joins,
        only the right for right joins.  Full outer joins are left
        unchanged.
\end{itemize}

\subsubsection*{Correctness}

Each case preserves the set of output rows:

\begin{itemize}[leftmargin=1.5em]
  \item \emph{Filter + SelectNames:} projection is a pure column subset;
        row identity is unchanged, so filter commutes.
  \item \emph{Filter + WithColumn:} the derived column is not in the
        predicate, so the boolean mask is identical whether computed before
        or after the column is added.
  \item \emph{Filter merge:} $\text{filter}(p_1)(\text{filter}(p_2)(t)) =
        \text{filter}(p_1 \wedge p_2)(t)$ by standard set-comprehension
        rules.
  \item \emph{Filter + Sort:} sorting is a permutation; filtering a
        permutation gives the same set of rows as filtering first and then
        permuting.
  \item \emph{Filter + Join on key columns (inner):} for an inner join the
        predicate on join keys can be applied to both sides before
        materialising the cross-product, because any row excluded by the
        predicate would not have contributed a matched pair.
\end{itemize}

\subsubsection*{Example}

\begin{lstlisting}
// Original plan
scan_parquet("employees.parquet")
  .sort({"salary"}, {true})
  .filter(col("dept") == lit("Engineering"))

// After predicate pushdown (filter pushed below sort)
scan_parquet("employees.parquet")
  .filter(col("dept") == lit("Engineering"))
  .sort({"salary"}, {true})
\end{lstlisting}

\subsubsection*{Performance benefit}

Reduces the number of rows before expensive operations (sort is
$O(n \log n)$; join is $O(n + m)$ average for hash-join).  A filter that
eliminates 90\% of rows before a sort turns a million-row sort into a
100\,000-row sort.

%------------------------------------------------------------------
\subsection{Projection Pushdown}
%------------------------------------------------------------------

\subsubsection*{Transformation}

\texttt{ApplyProjectionPushdown(node)} inspects a \texttt{SelectNamesNode}
or \texttt{SelectExprsNode} and attempts to reduce the number of columns
read/carried at lower levels:

\begin{itemize}[leftmargin=1.5em]
  \item \textbf{SelectNames over Filter:} Computes the union of the
        projected column set and the columns referenced by the filter
        predicate; inserts a \texttt{SelectNamesNode} below the filter to
        read only those columns.  An idempotency guard avoids inserting a
        redundant projection if one already exists.

  \item \textbf{SelectNames over WithColumn:} If the derived column is not
        needed in the output, the entire \texttt{with\_column} node is
        elided.  If it is needed, a narrowing projection is inserted below
        \texttt{with\_column} carrying only the columns required to compute
        the derived column plus the rest of the select list.

  \item \textbf{SelectExprs over Filter:} Analogous to the SelectNames
        case; the needed column set is the union of all columns referenced
        in the expression list and the filter predicate.
\end{itemize}

\subsubsection*{Correctness}

Projection pushdown only removes columns that are not consumed by any
downstream node.  The union with filter/expression column sets ensures that
no column that would affect computation is dropped.  The idempotency guard
(\texttt{IsExactSelectNames}) prevents infinite insertion of identical
projection nodes.

\subsubsection*{Example}

\begin{lstlisting}
// Wide table: 20 columns; we only need "name" and "salary"
scan_csv("employees.csv")            // 20 columns
  .filter(col("dept") == lit("Eng")) // references "dept"
  .select({"name", "salary"})

// After projection pushdown
scan_csv("employees.csv")            // 20 columns
  .select({"dept", "name", "salary"}) // <-- pushed down
  .filter(col("dept") == lit("Eng"))
  .select({"name", "salary"})
\end{lstlisting}

\subsubsection*{Performance benefit}

For wide Parquet files, reading only the required columns avoids
deserialising and allocating memory for unused column data.  For CSV the
benefit is primarily reduced memory allocation and fewer array iterations
in subsequent kernels.

%------------------------------------------------------------------
\subsection{Limit Pushdown}
%------------------------------------------------------------------

\subsubsection*{Transformation}

\texttt{ApplyLimitPushdown(node)} handles a \texttt{HeadNode} and pushes
the row limit earlier in the plan:

\begin{itemize}[leftmargin=1.5em]
  \item \textbf{Head over Head:} Two consecutive limits are merged into
        a single \texttt{HeadNode} with $n = \min(n_1, n_2)$.

  \item \textbf{Head over SelectNames / SelectExprs:} The head is pushed
        below the projection: the limit is applied before column projection,
        reducing the number of rows on which expressions are evaluated.

  \item \textbf{Head over WithColumn:} The limit is pushed below the
        column derivation: fewer rows need the derived expression evaluated.
\end{itemize}

\subsubsection*{Correctness}

\emph{Head merge:} $\text{head}(n_1)(\text{head}(n_2)(t)) =
\text{head}(\min(n_1,n_2))(t)$, since both limits take only the first $k$
rows in the same order.

\emph{Head + SelectNames/SelectExprs/WithColumn:} Column projections are
pure transformations that do not reorder or drop rows.  Therefore, applying
the row limit before or after the projection yields the same rows; it is
always cheaper to limit first.

\subsubsection*{Example}

\begin{lstlisting}
// Original
scan_csv("large.csv")
  .with_column("tax", col("salary") * lit(0.3))
  .head(10)

// After limit pushdown
scan_csv("large.csv")
  .head(10)                        // <-- pushed down
  .with_column("tax", col("salary") * lit(0.3))
\end{lstlisting}

\subsubsection*{Performance benefit}

Avoids computing the derived column (or expression projection) for rows
that will be discarded.  For a 1\,000\,000-row table with
\texttt{head(10)}, the \texttt{with\_column} evaluates only 10 rows
instead of 1\,000\,000.

%------------------------------------------------------------------
\subsection{Predicate Pushdown through GroupByAggregate}
%------------------------------------------------------------------

\subsubsection*{Transformation}

When a \texttt{FilterNode} sits directly above a
\texttt{GroupByAggregateNode} and every column referenced by the predicate
is a group-key column, the filter is pushed below the aggregate:

\[
\texttt{group\_by}(K)\texttt{.aggregate}(A)\texttt{.filter}(p)
\;\longrightarrow\;
\texttt{filter}(p)\texttt{.group\_by}(K)\texttt{.aggregate}(A)
\]

The check is performed by \texttt{CollectColumnsFromExpr} on the predicate
followed by a membership test against the key list.  If any predicate
column is not a group key (i.e.\ it is an aggregated output column such as
\texttt{salary\_mean}), the transformation is skipped to preserve
correctness.

\subsubsection*{Correctness}

A group-key column has the same value for every row in the group and
survives into the aggregate output unchanged.  Therefore a predicate that
only tests key columns produces the same boolean result whether it is
evaluated on the individual input rows or on the single representative row
emitted per group.  More formally: a group $g$ with key value $k$ is
included in the output iff $p(k)$ is true; filtering input rows with $p(k)$
before grouping eliminates exactly those groups, so the set of surviving
groups and their aggregate values are identical.  Predicates that reference
aggregated columns (e.g.\ \texttt{salary\_mean > 80000}) cannot be pushed
down because those values do not exist until after the aggregation
completes.

\subsubsection*{Example}

\begin{lstlisting}
// Original plan (filter after aggregation)
scan_csv("employees.csv")
  .join(scan_csv("departments.csv"), {"dept_id"}, "inner")
  .group_by({"dept_name"})
  .aggregate({{"salary", "mean"}})
  .filter(col("dept_name") == lit("Engineering"))

// After predicate pushdown through GroupByAggregate
scan_csv("employees.csv")
  .join(scan_csv("departments.csv"), {"dept_id"}, "inner")
  .filter(col("dept_name") == lit("Engineering"))  // <-- pushed down
  .group_by({"dept_name"})
  .aggregate({{"salary", "mean"}})
\end{lstlisting}

\subsubsection*{Performance benefit}

Rows belonging to excluded groups never enter the hash-map construction
phase of group\_by.  In the example above, only the Engineering rows
($\approx\!40\%$ of the table) are hashed; HR and Finance rows are
discarded cheaply by the filter scan.  The saving scales with the
selectivity of the key predicate and the cost of the aggregation kernel.

%------------------------------------------------------------------
\subsection{Dead WithColumn Elimination}
%------------------------------------------------------------------

\subsubsection*{Transformation}

\texttt{ApplyDeadWithColumnElimination} removes a \texttt{WithColumnNode}
whose derived column is provably never consumed by the immediately
downstream node.  Two structural patterns are handled:

\begin{itemize}[leftmargin=1.5em]
  \item \textbf{SelectExprs above WithColumn:} The set of column names
        referenced across all output expressions is collected via
        \texttt{CollectColumnsFromExpr}.  If the derived column name does
        not appear in that set, the \texttt{with\_column} node is bypassed:
        \texttt{select\_exprs(E).with\_column(c,e)} $\to$
        \texttt{select\_exprs(E)} with the input rewired to
        \texttt{with\_column}'s own input.

  \item \textbf{GroupByAggregate above WithColumn:} If the derived column
        is neither a group key nor referenced in any aggregation expression,
        the \texttt{with\_column} node is bypassed similarly.
\end{itemize}

The \texttt{SelectNames\,+\,WithColumn} case (projection of a named column
subset that does not include the derived column) is already handled by
\texttt{ApplyProjectionPushdown}; this pass covers the complementary
expression-select and aggregation cases.

\subsubsection*{Correctness}

A \texttt{with\_column} operation appends or replaces exactly one column
and leaves all other columns and rows unchanged.  If no downstream
expression references the derived column, the output of the plan is
identical whether or not the column was ever computed.  The membership
check against \texttt{CollectColumnsFromExpr} is conservative: it scans
the entire expression tree of every downstream expression, so no
reference can be missed.

\subsubsection*{Example}

\begin{lstlisting}
// Original
scan_csv("employees.csv")
  .with_column("double_salary", col("salary") * lit(2))
  .select(std::vector<Expr>{col("id").alias("eid"), col("name")})
  // "double_salary" not referenced in select expressions

// After Dead WithColumn Elimination
scan_csv("employees.csv")
  .select(std::vector<Expr>{col("id").alias("eid"), col("name")})
\end{lstlisting}

In the smoke test this fires together with Constant Filter Elimination in a
single optimiser pass, reducing a four-node logical plan to two nodes:

\begin{center}
\texttt{scan \textrightarrow{} filter(1==1) \textrightarrow{} with\_column("double\_salary") \textrightarrow{} select\_exprs}
\\[0.3em]
$\Downarrow$
\\[0.3em]
\texttt{scan \textrightarrow{} select\_exprs}
\end{center}

\subsubsection*{Performance benefit}

Eliminates the expression evaluation kernel for the dead column entirely --
no intermediate array is allocated or populated.  For expressions involving
expensive string operations or long arithmetic chains this saving can be
substantial.  The optimisation also reduces plan depth, which shortens the
executor's dispatch loop.

%------------------------------------------------------------------
\subsection{Consecutive Projection Merging}
%------------------------------------------------------------------

\subsubsection*{Transformation}

\texttt{ApplyProjectionMerging} inspects a \texttt{SelectNamesNode} and
collapses it with a projection immediately below it:

\begin{itemize}[leftmargin=1.5em]
  \item \textbf{SelectNames above SelectNames:} The outer (downstream)
        projection is rewired directly to the inner node's input, discarding
        the inner \texttt{SelectNamesNode} entirely.  Because the outer
        column list is already a subset of the inner one (if the plan is
        valid), no column safety check is required.

  \item \textbf{SelectNames above SelectExprs:} The expressions whose
        output names match the outer column list are filtered and reordered
        to match the outer column order, producing a single
        \texttt{SelectExprsNode}.  If any required name cannot be resolved,
        the transformation is skipped.
\end{itemize}

\subsubsection*{Correctness}

A projection is a pure column-subsetting operation.  Applying two
projections $\pi_A(\pi_B(R))$ where $A \subseteq B$ is equivalent to
$\pi_A(R)$ by the standard projection absorption law.  For the
\texttt{SelectNames\,+\,SelectExprs} case, only the expressions whose
declared output names are in the outer set are retained; all others are
unreachable and their removal cannot affect the result.

\subsubsection*{Example}

\begin{lstlisting}
// Original (redundant wide projection followed by narrow one)
scan_csv("employees.csv")
  .select({"id", "name", "dept_id", "age", "salary"})
  .select({"id", "name", "salary"})

// After consecutive projection merging
scan_csv("employees.csv")
  .select({"id", "name", "salary"})
\end{lstlisting}

\subsubsection*{Performance benefit}

Eliminates an entire projection pass -- no intermediate column-subset array
is allocated.  More importantly it reduces plan depth, which shortens the
dispatch loop in \texttt{ExecutePlan} and creates a shallower structure for
other pushdown rules to traverse.

%------------------------------------------------------------------
\subsection{Join Predicate Splitting}
%------------------------------------------------------------------

\subsubsection*{Transformation}

\texttt{ApplyJoinPredicateSplitting} fires when a \texttt{FilterNode} sits
directly above an inner \texttt{JoinNode} and the predicate is a
conjunction.  The transformation proceeds as follows:

\begin{enumerate}[leftmargin=1.5em]
  \item \texttt{SplitConjunction} recursively decomposes the predicate on
        \texttt{AND} boundaries into a flat list of clauses.
  \item \texttt{TryGetOutputColumns} is called on each join input to
        determine its output column set.  The function traverses
        \texttt{SelectNames}, \texttt{SelectExprs}, \texttt{Filter},
        \texttt{Sort}, \texttt{Head}, \texttt{WithColumn},
        \texttt{GroupByAggregate}, and \texttt{Join} nodes; it returns
        \texttt{nullopt} for bare scan nodes whose schema is unknown at
        plan time.
  \item Each clause is classified: non-key columns that belong exclusively
        to the left schema are pushed to the left input; those belonging
        exclusively to the right schema are pushed to the right input;
        clauses that span both sides (or reference only join keys) remain
        above the join.
  \item New \texttt{FilterNode}s are inserted into the respective inputs,
        and any remaining clauses are reassembled with
        \texttt{CombineClauses} above the rebuilt join.
\end{enumerate}

The optimisation applies only to inner joins; outer, left, and right joins
are left unchanged because null-padding semantics make independent
pre-filtering incorrect.

\subsubsection*{Correctness}

For an inner join, a row pair $(L, R)$ appears in the output iff the join
key matches \emph{and} all predicates evaluate to true.  If clause $p_L$
references only columns from $L$ and clause $p_R$ references only columns
from $R$, then:

\[
\sigma_{p_L \wedge p_R}(L \bowtie R) = (\sigma_{p_L} L) \bowtie (\sigma_{p_R} R)
\]

This follows because $p_L(L, R) = p_L(L)$ and $p_R(L, R) = p_R(R)$, so
filtering can be applied to each side independently before the join is
formed.  The key-only check avoids pushing predicates on join-key columns
(which would interact with the null-handling logic in the existing predicate
pushdown rule).

\subsubsection*{Example}

\begin{lstlisting}
// Original: combined filter above join
scan_csv("employees.csv").select({"id","name","dept_id","age","salary"})
  .join(
    scan_csv("departments.csv").select({"dept_id","dept_name","location"}),
    {"dept_id"}, "inner")
  .filter(col("age") > lit(30) & col("dept_name").starts_with("E"))

// After join predicate splitting (age -> left, dept_name -> right)
scan_csv("employees.csv").select({...}).filter(col("age") > lit(30))
  .join(
    scan_csv("departments.csv").select({...}).filter(col("dept_name").starts_with("E")),
    {"dept_id"}, "inner")
\end{lstlisting}

\subsubsection*{Performance benefit}

Both inputs to the join are filtered before the hash index is built and
probed.  In the example, only Engineering rows survive on the right (1 of
3 departments), and only employees older than 30 survive on the left (5 of
10).  The hash index is built over 1 row instead of 3, and only 5 left rows
are probed instead of 10 -- a combined reduction of roughly $5\times$ in
join work on this dataset.  The saving scales with selectivity and is
particularly large when one side is highly selective.

%------------------------------------------------------------------
\subsection{Filter Merge (Predicate Conjunction)}
%------------------------------------------------------------------

This sub-transformation is embedded inside \texttt{ApplyPredicatePushdown}
(the \emph{Filter over Filter} case).  Two consecutive
\texttt{FilterNode}s are collapsed into one whose predicate is the AND of
both.  This reduces plan depth and allows subsequent passes to push the
combined predicate further down.

%------------------------------------------------------------------
\subsection{Multi-pass Convergence Strategy}
%------------------------------------------------------------------

The optimizer applies \texttt{OptimizeTreeOnce} (a bottom-up traversal
followed by one application of all rule functions per node) in a loop of
at most 8 passes.  After each pass, a string fingerprint of the plan tree
is computed; if the fingerprint does not change the loop exits early.  This
handles chains of transformations where one pass creates an opportunity
exploited by the next (e.g.\ predicate pushdown enabling projection
pushdown).

%--------------------------------------------------------------------------
\section{Memory Management}
%--------------------------------------------------------------------------

The library uses the following C++ resource-management idioms throughout:

\begin{itemize}[leftmargin=1.5em]
  \item \textbf{RAII via \texttt{shared\_ptr}:}  All heap-allocated
        objects (\texttt{ExprNode}, \texttt{PlanNode},
        \texttt{arrow::Array}, \texttt{arrow::Schema}) are owned through
        \texttt{shared\_ptr}.  Reference-counted ownership ensures prompt
        cleanup as DataFrames and plan nodes go out of scope.

  \item \textbf{Move semantics:}  \texttt{EagerDataFrame} holds its column
        vector and schema by value; factory functions and transformation
        methods use \texttt{std::move} to transfer ownership without
        unnecessary copying.  Column vectors passed into
        \texttt{from\_ordered\_columns} are moved.

  \item \textbf{Value types for DataFrames:}  Both \texttt{EagerDataFrame}
        and \texttt{LazyDataFrame} are copyable/movable value types;
        copying a \texttt{LazyDataFrame} is cheap (it copies a
        \texttt{shared\_ptr}).

  \item \textbf{Arrow memory pool:}  All Arrow array builders use the
        default Arrow memory pool, giving consistent allocation tracking.
\end{itemize}

No raw \texttt{new} / \texttt{delete} calls appear in the library code.
Valgrind-clean operation follows from Arrow's own memory management.

%--------------------------------------------------------------------------
\section{Build Instructions}
%--------------------------------------------------------------------------

\subsection{Prerequisites}

\begin{itemize}[leftmargin=1.5em]
  \item CMake $\ge$ 3.20
  \item A C++20-capable compiler (GCC $\ge$ 11 or Clang $\ge$ 13)
  \item Apache Arrow $\ge$ 12 (with Parquet support) -- install via
        conda, vcpkg, or from source.
  \item Graphviz (\texttt{dot}) -- required at runtime for
        \texttt{explain()}.
\end{itemize}

\subsection{Building}

\begin{verbatim}
cmake -B build -DCMAKE_BUILD_TYPE=Release
cmake --build build -j$(nproc)
\end{verbatim}

The build produces:
\begin{itemize}[leftmargin=1.5em]
  \item \texttt{build/libdataframelib.so} (or \texttt{.dylib} on macOS)
  \item \texttt{build/csv\_smoketest} -- a smoke-test executable
\end{itemize}

\subsection{Running the smoke test}

\begin{verbatim}
./build/csv_smoketest
\end{verbatim}

\subsection{Installation}

\begin{verbatim}
cmake --install build --prefix /usr/local
\end{verbatim}

This installs the shared library to \texttt{lib/} and the headers to
\texttt{include/dataframelib/}.

%--------------------------------------------------------------------------
\section{Example Usage}
%--------------------------------------------------------------------------

\subsection{Eager Mode}

\begin{lstlisting}
#include "dataframelib/dataframelib.h"
using namespace dataframelib;

auto df = EagerDataFrame::read_parquet("data.parquet");
auto result = df.filter(col("age") > lit(30))
                .select({"name", "salary"});
result.write_csv("output.csv");
\end{lstlisting}

\subsection{Lazy Mode with Optimizer}

\begin{lstlisting}
auto df = LazyDataFrame::scan_parquet("data.parquet");
auto result = df
    .filter(col("age") > lit(30))
    .select({"name", "salary"})
    .group_by({"dept"})
    .aggregate({{"salary", "mean"}});

result.explain("plan.png");   // renders logical + optimized DAGs
auto output = result.collect(); // triggers optimisation + execution
output.write_csv("summary.csv");
\end{lstlisting}

\end{document}
